
\chapter{Chapter 3. Datasets}

This chapter describes the two datasets used in this thesis: the multi-source training dataset used to fine-tune the base embedding model, and the evaluation benchmark used to measure retrieval performance. These two datasets are kept separate so that there is no data leakage. This means no training example appears in any evaluation set. The training data comes from question answering, paraphrasing, and natural language inference corpora that are open to the public. The NanoBEIR benchmark is used exclusively for evaluation.

\section{Training Data}

The pre-trained model is fine-tuned on a mixture of seven publicly available datasets sourced from the Hugging Face Hub. The seven sources come from different domains and pair types, which helps reduce the risk of overfitting and improve generalization. Each dataset has either (query, positive passage) pairs or (query, positive passage, negative passage) triplets that will be handled by Multiple Negatives Ranking Loss (MNRL) during training.

\subsection{Dataset Sources}

The seven training datasets cover five broad semantic similarity tasks: passage retrieval, title to abstract prediction, duplicate question detection, natural language inference, and open-domain question answering. Table \ref{tab:training_data} gives an overview of each dataset, including its pair type, raw size, sampling weight, and the number of training examples that are taken from it during each training period, which is described in Section 3.1.2.

\begin{table}[H]
  \centering
  \caption[Training dataset summary.]{Training dataset summary. Effective rows show the contribution of each dataset to the training mixture after weighted batch sampling. Total rows from all sources sum to 58.73M; effective rows sum to about 4.19M.}
  \label{tab:training_data}
  \resizebox{\columnwidth}{!}{%
  \begin{tabular}{| c | l | l | r | r | r |}
    \hline
    \textbf{S.No.} & \textbf{Dataset} & \textbf{Pair Type} & \textbf{Total Rows} & \textbf{Weight} & \textbf{Effective Rows} \\ [0.5ex] \hline
    1 & msmacro\_triplet          & Query to Passage (triplet)         & 13.64M & 0.15 & $\sim$1.94M \\ \hline
    2 & s2orc\_title\_abstract    & Title to Abstract (pair)           & 41.77M & 0.01 & $\sim$0.40M \\ \hline
    3 & quora\_dup\_triplet       & Duplicate question (triplet)       & 2.79M  & 0.15 & $\sim$0.40M \\ \hline
    4 & nli\_for\_simcse\_triplet & Entailment/Contradiction (triplet) & 0.27M  & 2    & $\sim$0.52M \\ \hline
    5 & natural\_questions        & Question to Answer (pair)          & 0.10M  & 4    & $\sim$0.38M \\ \hline
    6 & trivia\_qa\_triplet       & Question to Answer (triplet)       & 0.06M  & 4    & $\sim$0.23M \\ \hline
    7 & hotpotqa                  & Multi-hop Q to A (triplet)         & 0.08M  & 4    & $\sim$0.32M \\ \hline
    \multicolumn{3}{| l |}{\textbf{Total}} & \textbf{58.73M} & & \textbf{$\sim$4.19M} \\ \hline
  \end{tabular}}
\end{table}

\textbf{MS MARCO Triplets} (\textit{msmacro\_triplet}) is derived from the Microsoft MAchine Reading COmprehension dataset \cite{bajaj2016ms}, which has 13.64 million (query, passage, negative passage) triplets and is the biggest source in the mix. Queries are real Bing search queries that people have used, and passages are taken from web pages.

\textbf{Semantic Scholar Open Research Corpus} (\textit{s2orc\_title\_abstract}) provides 41.77 million (title, abstract) pairs taken from scientific publications in a wide range of academic fields \cite{lo2020s2orc}. This dataset trains a model to connect document titles with their abstracts, which works as an equivalent for retrieving documents.

\textbf{Quora Duplicate Triplets} (\textit{quora\_dup\_triplet}) contains 2.79 million (question, duplicate question, non-duplicate question) triplets drawn from the Quora Question Pairs dataset. This source trains the model to put questions that are semantically similar close to each other and helps with retrieval tasks where queries may come in paraphrased form.

\textbf{NLI for SimCSE Triplets} (\textit{nli\_for\_simcse\_triplet}) adapts natural language inference corpora following the SimCSE training recipe \cite{gao2021simcse}, resulting in 0.27 million triplets (premise, entailment hypothesis, contradiction hypothesis). You can think of entailment pairs as strong positives and contradiction pairs as strong negatives.

\textbf{Natural Questions} (\textit{natural\_questions}) contains 0.10 million (question, answer passage) pairs taken from Google's Natural Questions dataset \cite{kwiatkowski2019natural}, in which questions are real user queries submitted to the Google search engine, and answers are Wikipedia passages marked by human annotators.

\textbf{TriviaQA Triplets} (\textit{trivia\_qa\_triplet}) provides 0.06 million (trivia question, supporting passage, negative passage) triplets from the TriviaQA reading comprehension dataset \cite{joshi2017triviaqa}. Questions require a lot of factual knowledge and come with passages from Wikipedia and other web documents that back them up.

\textbf{HotpotQA} (\textit{hotpotqa}) adds 0.08 million multi-hop triplets (question, supporting passage, negative passage) \cite{yang2018hotpotqa}. To answer HotpotQA questions, you have to look at more than one document and think about what it means. This is a more difficult retrieval task than answering a single-hop question.

\subsection{Weighted Batch Sampling}

Rather than concatenating all seven datasets and sampling uniformly, the training pipeline uses weighted batch sampling, where each dataset is given a scalar weight that controls how many examples are expected to be taken from it in each training step. Let $w_i$ denote the weight of dataset $i$ and $N_i$ its total number of examples. The effective contribution of dataset $i$ to the training mixture is proportional to $w_i \cdot N_i$. This is scaled so that the total contributions from all datasets add up to about 4.19 million useful examples (see Table \ref{tab:training_data}).

This is useful in two ways. First, it stops the two biggest sources, MS MARCO and S2ORC, from flooding the gradient signal with a single retrieval pattern. Second, it allows smaller domain-specific question-answering datasets like Natural Questions, TriviaQA, and HotpotQA to contribute meaningfully to training despite their limited raw size. The effective training mixture of about 4.19 million examples per epoch is much smaller than the original 58.73 million examples. This cuts down on training time while keeping the diversity of the training examples.

\section{Evaluation Benchmark: NanoBEIR}

All experiments in this thesis are evaluated on the NanoBEIR benchmark. NanoBEIR is based on the full BEIR benchmark, which was considered but not used because it would have been too expensive to encode corpora with millions of passages for each compression version being studied.

NanoBEIR \cite{nanobeir2024} is a lightweight subset of the Benchmarking Information Retrieval (BEIR) benchmark \cite{thakur2021beir}. It was made by sampling each subset in the BEIR dataset with 50 queries and a small corpus to go with it. As part of the \textit{sentence-transformers} ecosystem, it is shared on the Hugging Face Hub. Its purpose is to allow quick model comparisons without the computational cost of storing the full BEIR corpora, which makes it ideal for iterative development and ablation studies.

Table \ref{tab:nanobeir} shows a summary of all 13 NanoBEIR subsets that are used in this thesis. The subsets cover a wide range of retrieval tasks, such as fact checking, answering questions, entity retrieval, citation prediction, biomedical search, and argument retrieval. They also cover a wide range of domains, from general web material to financial documents to scientific literature. There are a total of 649 queries and 57,390 corpus documents across all 13 subsets. Each subset has 50 queries (49 in the case of NanoTouche2020) and a corpus of between 2,260 and 6,100 documents.

\begin{table}[H]
  \centering
  \caption[NanoBEIR benchmark subsets.]{NanoBEIR benchmark subsets used for evaluation. Except for NanoTouche2020, which has 49 queries, all subsets have 50 queries. It can have anywhere from 2,260 to 6,100 documents. Altogether, there are 649 queries and 57,390 corpus records.}
  \label{tab:nanobeir}
  \resizebox{\columnwidth}{!}{%
  \begin{tabular}{| c | l | r | r | l |}
    \hline
    \textbf{S.No.} & \textbf{Subset} & \textbf{Queries} & \textbf{Corpus} & \textbf{Task (Query Type to Document Type)} \\ [0.5ex] \hline
    1  & NanoArguAna         & 50 & 3,690 & Counter-argument retrieval (Claim to Argument)          \\ \hline
    2  & NanoClimateFEVER    & 50 & 3,460 & Fact-checking (Claim to Document)                       \\ \hline
    3  & NanoDBPedia         & 50 & 6,100 & Entity retrieval (Query to Document)                    \\ \hline
    4  & NanoFEVER           & 50 & 5,050 & Fact-checking (Claim to Document)                       \\ \hline
    5  & NanoFiQA2018        & 50 & 4,650 & Financial question answering (Question to Document)     \\ \hline
    6  & NanoHotpotQA        & 50 & 5,140 & Multi-hop question answering (Question to Document)     \\ \hline
    7  & NanoMSMARCO         & 50 & 5,090 & Web search / general QA (Question to Passage)           \\ \hline
    8  & NanoNFCorpus        & 50 & 3,000 & Biomedical information retrieval (Query to Document)    \\ \hline
    9  & NanoNQ              & 50 & 5,090 & Open-domain question answering (Question to Document)   \\ \hline
    10 & NanoQuoraRetrieval  & 50 & 5,100 & Duplicate detection (Question to Question)              \\ \hline
    11 & NanoSCIDOCS         & 50 & 2,260 & Citation prediction (Title to Document)                 \\ \hline
    12 & NanoSciFact         & 50 & 2,970 & Scientific fact-checking (Claim to Document)            \\ \hline
    13 & NanoTouche2020      & 49 & 5,790 & Conversational argument retrieval (Question to Argument)\\ \hline
    \multicolumn{2}{| l |}{\textbf{Total}} & \textbf{649} & \textbf{57,390} & \\ \hline
  \end{tabular}}
\end{table}

The diversity of NanoBEIR tasks is particularly important for evaluating compression methods. For a standard benchmark, a good average score might hide that the method fails for some other type of task, but with NanoBEIR, we test the compression method on 13 different subsets.

